\documentclass{article}
\usepackage[margin=50pt]{geometry}
\usepackage[fontsize=15pt]{fontsize}
\usepackage{lineno}
\linenumbers
% graphics packages
\usepackage{graphicx}
\usepackage[obeyclassoptions,mode=tex]{standalone}
\usepackage{tikz}
\usetikzlibrary{backgrounds}
\usetikzlibrary{shapes.geometric}
\usetikzlibrary{positioning}
\usetikzlibrary{automata}
\usetikzlibrary{tikzmark}
\usetikzlibrary{patterns}
\usetikzlibrary{arrows}
\tikzset{every state/.style={minimum size=1pt}}
\usepackage{tikz-cd}

% Maths
\usepackage{amssymb,stmaryrd,amsthm,amsmath}

% links inside the document
\usepackage[electronic,hyperref,xcolor,cleveref]{knowledge}
\knowledgeconfigure{notion}

%Theorem environments
\newtheorem{theorem}{Theorem}
\newtheorem{definition}[theorem]{Definition}
\newtheorem{lemma}[theorem]{Lemma}
\newtheorem{corollary}[theorem]{Corollary}
\newtheorem{example}[theorem]{Example}
\newtheorem{remark}[theorem]{Remark}
\newtheorem{conjecture}[theorem]{Conjecture}
\newtheorem{convention}[theorem]{Convention}
% Tables 
\usepackage{booktabs}
\usepackage{varwidth}

% Packages for macro definitions
\usepackage{xparse}
\usepackage{xpatch}
\usepackage{tokcycle}
\usepackage{ifthen}



% Proofs
\usepackage{bussproofs}

% Colors 
\definecolor{KLA4}{RGB}{0,78,125}
\definecolor{KLB1}{RGB}{49,62,72}
\knowledgestyle{intro notion}{color={KLA4}, emphasize}
\knowledgestyle{notion}{color={KLB1}}

% Algorithms
\usepackage{algorithm2e}
\Crefname{algocfline}{Algorithm}{Algorithms}
\crefname{algocfline}{Algorithm}{Algorithms}
\crefname{algocf}{Algorithm}{Algorithms}
\Crefname{algocf}{Algorithm}{Algorithms}


% we include whatever the user wants to include in the header
\graphicspath{ {./tree_amal/} }

% we include libraries (tex files) usually written in the `lib` directory
\input{lib/aliaume}
\input{lib/arka}
\input{lib/maths}
\input{lib/knowledges.kl}

% sample knowledge
\knowledge{notion}
 | kl-usage

% Knowledge logo
\newcommand{\klogo}{%
\begin{tikzpicture}[scale=0.2,line/.style={draw, line width=0.2pt, line cap=round, line join=round}]
\coordinate (A00) at (0,0);
\coordinate (A01) at (0,1);
\coordinate (A10) at (1,0);
\coordinate (B10) at (1,0.2);
\coordinate (B01) at (0.2,1);

\coordinate (C01) at (0.4,0.7);
\coordinate (C10) at (0.7,0.4);
\coordinate (C12) at (0.4,1.2);
\coordinate (C21) at (1.2, 0.4);
\coordinate (C22) at (1.2, 1.2);

\coordinate (D00) at (C10);
\coordinate (D01) at (0.8,0.5);
\coordinate (D10) at (0.8,0.3);

\coordinate (E01) at (0.3,0.7);
\coordinate (E10) at (0.5,0.7);

\draw[line] (B01) -- (A01) -- (A00) -- (A10) -- (B10);
\draw[line] (C01) -- (C12) -- (C22) -- (C21) -- (C10);

\draw[line] (D01) -- (D00) -- (D10);
\draw[line] (E01) -- (E10);

\end{tikzpicture}%
}
 
\title{Computability of Equivariant Gröbner bases}

\author{Arka Ghosh and Aliaume Lopez}

\begin{document}
% Generate the title page
\maketitle
% Print the abstract
\begin{abstract}
    Let \(\mathbb{K}\) be a field, \(\mathcal{X}\) be an infinite set (of indeterminates), and \(\mathcal{G}\) be a group acting on \(\mathcal{X}\). An ideal in the polynomial ring \(\mathbb{K}[\mathcal{X}]\) is called equivariant if it is invariant under the action of \(\mathcal{G}\). We show Gröbner bases for equivariant ideals are computable are hence the equivariant ideal membership is decidable when \(\mathcal{G}\) and \(\mathcal{X}\) satisfies the Hilbert's basis property, that is, when every equivariant ideal in \(\mathbb{K}[\mathcal{X}]\) is finitely generated. Moreover, we give a sufficient condition for the undecidability of the equivariant ideal membership problem. This condition is satisfied by the most common examples not satisfying the Hilbert's basis property. Our results imply decidability of solvability of orbit-finite systems of linear equations and the reachability problem for reversible data Petri nets for a large class of data domains.
\end{abstract}

\noindent
\klogo\ This document uses \href{https://ctan.org/pkg/knowledge}{knowledge}:
a \kl[kl-usage]{notion} points to its \intro[kl-usage]{definition}. 


% Include the content of the paper
%!TEX root = ../atomic_jsc.tex
% LTeX: language=en
%
\section{Introduction}
%
Let $\K$ be a field and $\Indets$ be a countable (i.e., finite or countably infinite) set of variables.
Let $\monoid$ be a monoid acting injectively on $\Indets$
(i.e., the function $x\mapsto \gelem\cdot x$ is injective for every $\gelem\in\monoid$).
An ideal in $\poly{\K}{\Indets}$ is called $\monoid$-equivariant if it is invariant under the action of $\monoid$.
For instance, let $\inc{\N}$ be the monoid of strictly monotonic functions from $\N$ to itself and the action on $\Indets = \{x_1,x_2,\dots\}$ is defined as $\gelem \cdot x_i = x_{\gelem(i)}$ for $\gelem \in \inc{\N}$ and $x_i\in \Indets$.
Then the ideal $\idl[S]_0$ of polynomials in $\poly{\K}{\Indets}$ whose coefficients add up to zero is a $\inc{\N}$-equivariant ideal.
In \cite{COHEN67} Cohen showed that the classical Hilbert's basis theorem \arka{cite wiki} extends to $\inc{\N}$-equivariant ideals: every $\inc{\N}$-equivariant ideal in $\poly{\K}{\Indets}$ is generated by an orbit-finite set of polynomials.
For instance, $\idl[S]_0$ is generated by the set $\setof{x_i - 1}{i\in\N}$,
which forms a single orbit under the action of $\inc{\N}$.
Let $\symgr{\Indets}$ denote the group of all permutations of $\Indets$.
Since every $\symgr{\Indets}$-equivariant ideal is also $\inc{\N}$-equivariant,
as a corollary of Cohen's result one gets that $\symgr{\Indets}$-equivariant ideals are also finitely generated.
The results of Cohen were rediscovered in \cite{AH07} and generalised in \cite{HISU12,LauSno26,GHOLAS24}.

To explain these results we need to give a few definitions.
We say that the action $\monoid\actson\Indets$ is:
\begin{enumerate}
\item \intro{ordered} if it respects a total order $\prec$ on $\Indets$ (i.e., $\gelem \cdot x \prec \gelem \cdot y$ for every $x \prec y$ and $\gelem\in\monoid$),
\item \intro{well-ordered} if it respects a well-order on $\Indets$,
\item \intro{nice} if the the divisibility up to $\monoid$ relation ($\monelt[p] \gdivleq[\monoid] \monelt[q]$ if there exists $\pi\in\monoid$ such that $\pi\cdot\monelt[p]$ divides $\monelt[q]$) is a well-quasi-order on the set $\mon{\Indets}$ of monomials in $\Indets$,
\item \intro{nicely ordered} if it is \kl{nice} and \kl{ordered},
\item \intro{nicely well-ordered} if it is \kl{nice} and \kl{well-ordered}.
\end{enumerate}
%
Observe that an \kl{ordered} monoid action must be injective.
The action $\symgr{\Indets}\actson\Indets$ is \kl{nice} but not \kl{ordered}.
The action of $\inc{\N}$ mentioned above is \kl{nicely well-ordered}.
In \arka{refer sec} we discuss several examples of monoid actions and discuss which of the above properties they satisfy.


The article \cite{HISU12} showed proved $\monoid$-equivariant ideals are orbit-finitely generated when $\monoid\actson\Indets$ is \kl{nicely well-ordered},
and used it to prove the independent set conjecture in algebraic statistics.
This result was further generalised in \cite{GHOLAS24,LauSno26} which showed that:
\begin{enumerate}
\item $\monoid$-equivariant ideals are orbit-finitely generated when $\monoid\actson\Indets$ is \kl{nicely ordered} (i.e., they get rid of the well-foundedness condition on the ordering on variables).
\item Conversely, if $\monoid$-equivariant ideals are orbit-finitely generated then $\monoid\actson\Indets$ is \kl{nice}.
\end{enumerate}
%
It has been conjectured by Pouzet that if $\monoid\actson\Indets$ is \kl{nice},
then there exists another monoid $\mathcal{N}$ acting on $\Indets$ such that 
\begin{enumerate}
\item for every finite subset $S$ of $\Indets$ and $\gelem\in\mathcal{N}$,
there exists $\sigma\in\monoid$ such that $\sigma\cdot x = \gelem\cdot x$ for every $x\in S$.
\item $\mathcal{N}\actson\Indets$ is \kl{nicely ordered}. 
\end{enumerate}
%
Monoid actions $\monoid\actson\Indets$ satisfying these conditions will be called \intro{nicely orderable}.
%
\begin{remark}
It is tempting to conjecture that $\mathcal{N}$ can even be a sub-monoid of $\monoid$.
However this stronger conjecture would be false as there is no sub-monoid $\mathcal{N}$ of the group of finitary permutations of $\Indets$ (permutations $\pi$ that $\pi(x) \neq x$ only for finitely many $x\in\Indets$) that satisfy the two above conditions.
Observe however $\inc{\N}$ is a monoid that satisfies these conditions but is not a sub-monoid of the group of finitary permutations.
\end{remark}
%
\begin{example}
We give an example of a \kl{nicely ordered} action that is not well-ordered.
Let $\prec$ be a dense linear order on $\Indets$ without endpoints (such an order can be constructed from a bijection of $\Indets$ with the set of rational numbers), and $\monoid$ be the set of all strictly monotonic bijections from $\Indets$ to itself.
Clearly $\monoid\actson\Indets$ is not well-ordered.
However $\gdivleq[\monoid]$ is a well-quasi-order on $\mon{\Indets}$ due to Higman's lemma.
\arka{we discuss a proof scheme to prove niceness later..}
\end{example}
%
The result of \cite{HISU12} was extended in \cite{BRDR11} which defined Gr\"{o}bner bases for $\monoid$-equivariant ideals and showed that they are orbit-finite and computable by an appropriate extension of the classical Buchberger's 
\arka{gave 'a' definition. Mention nicely well-ordered case before. List the contributions}
algorithm in the equivariant setting when $\monoid\actson\Indets$ is \kl{nicely well-ordered} and satisfies certain computability conditions that are necessary for running the equivariant Buchberger's algorithm.
Our paper extends the results of \cite{GHOLAS24,LauSno26} and generalise the results of \cite{BRDR11} to the case when the action of $\monoid\actson\Indets$ is \kl{nicely ordered}.
%
\subsection{Contributions}
\arka{cite theorems and sections. Also list the contributions}

As our first main result we define Gr\"{o}bner bases for $\monoid$-equivariant ideals and show that when $\monoid\actson\Indets$ is \kl{nicely ordered} and satisfies certain computability criterions,
every $\monoid$-equivariant ideal has an orbit-finite Gr\"{o}bner basis and
there is an algorithm that takes an orbit-finite set $H\subseteq\poly{\K}{\Indets}$ as input and outputs an orbit-finite Gr\"{o}bner base for $\monoid$-equivariant ideal generated by $H$ (an orbit-finite set
$(\monoid\cdot f_1)\cup\dots\cup(\monoid\cdot f_n)$ is represented by the set $\set{f_1,\dots,f_n}$).

As a corollary of our results we get that the $\monoid$-equivariant ideal membership problem is decidable.
Moreover, given orbit-finite Gr\"{o}bner bases for $\monoid$-equivariant ideals $\idl[I]$ and $\idl[J]$ one can compute orbit-finite Gr\"{o}bner bases for $\monoid$-equivariant ideals generated by $\idl[I]\cup\idl[J]$, $\idl[I]\cap\idl[J]$ and $\idl[I]\idl[J]$.

We complement our positive results by giving a sufficient condition on $\monoid\actson\Indets$ for the undecidability of the $\monoid$-equivariant ideal membership problem.
This sufficient condition is satisfied by most common actions $\monoid\actson\Indets$ which are not \kl{nice} (and hence there are $\monoid$-equivariant ideals that are not finitely generated). 

Finally, we sketch some applications of our results for solving the membership problem for orbit-finitely spanned vector spaces that has received the attention of the theoretical computer science community in recent years.
\arka{cite papers}
%
\subsection{Comparison with previous work}
%
We quickly describe how we improve the state of the art on Gr\"{o}bner bases by comparing our results to our immediate predecessors \cite{GHOLAS24,BRDR11}.

In addition to proving orbit-finite generation of $\monoid$-equivariant ideals \cite{GHOLAS24} also show that for some specific \kl{nicely ordered} actions $\monoid\actson\Indets$,
one can reduce the $\monoid$-equivariant ideal membership problem to the $\mathcal{N}$-equivariant ideal membership problem for some nicely well-ordered action $\mathcal{N}\actson\Indets$
(we point the reader to \arka{refer} for more details).
However, it is unclear whether their technique can be generalised to the class of all \kl{nicely ordered} actions,
at least the authors of \cite{GHOLAS24} do not claim that to be the case.
In particular, it seems difficult to extend the technique of \cite{GHOLAS24} to actions $\monoid\actson\Indets$ where $\Indets$ is the structure of a densely coloured dense linear order or the ordered dense meet tree and $\monoid$ is the group of automorphisms of the structure.\arka{refer to def}
Both of these actions are \kl{nicely ordered}.
 
Furthermore, the reduction to the nicely well-ordered actions only solves the ideal membership problem.
Extending the notion of Gr\"{o}bner base to the general setting of \kl{nicely ordered} actions allows us to also show effectivity of simple operations as sum and intersection of equivariant ideals,
which would not be possible using the results of \cite{GHOLAS24}.
At the same time it opens up the opportunity to extend complicated operations such as computing the primary decomposition and radical of ideals to the equivariant setting. 

Let us now explain the difficulty in generalising the concepts of \kl{Gr\"{o}bner bases} and \kl{Buchberger's algorithm} from \kl{nicely well-ordered} actions to \kl{nicely ordered actions}.
When the action $\monoid\actson\Indets$ preserves an well-order $\prec$ on $\Indets$,
one can define a well-order on $\mon{\Indets}$ which is compatible with the monoid action and divisibility,
by extending $\prec$ to $\mon{\Indets}$ with some suitable ordering such as co-lex or grevlex \arka{refer to defintions}.
However the same is not possible for all \kl{nicely ordered} actions.
%In the case when the action $\monoid\actson\Indets$ is \kl{nicely ordered} and neither the order $\prec$ preserved by $\monoid\actson\Indets$ nor its reverse is well-founded,
%it is not possible to define a well-order on $\Indets$ which is compatible with the action as any such order must agree with $\prec$ or its reverse
%(for instance co-lex agrees with $\prec$ and lex agrees with its reverse
%\arka{this is not always true. But it is true for ordered atoms. Change later as well}).
This means that sequence of reductions that merely decreases a ordering of monomials based on $\prec$ may not terminate.
%
\begin{example}
For example, consider the action $\inc{\OrderAtoms}\actson\OrderAtoms$.
Let $\monlt$ be a monomial ordering on $\mon{\OrderAtoms}$ that is $\monoid$-equivariant.
Pick $x < y \in \OrderAtoms$.
If $x \monlt y$, then $x' < y'$ for every $x' < y'\in \OrderAtoms$, because for every such $x'$ and $y'$ there exists some $\pi\in\inc{\OrderAtoms}$ such that $\pi\cdot x = x'$ and $\pi\cdot y = y'$.
But then $\monlt$ agrees with $<$ and hence cannot be a well-order.
Similarly, if $x \monlt y$ then $\monlt$ agrees with the reverse of $\prec$ and hence cannot be a well-order.


Assume $(\Indets,\prec)$ is a dense linear order without endpoints.
If we try to reduce any polynomial $x\in\poly{\K}{\Indets}$ using the set of polynomials $\setof{y - z}{y \prec z}$ we get an infinite sequence
$x \xrightarrow{} x_1 \xrightarrow{} x_2 \xrightarrow{}\dots$
for some $x \succ x_1 \succ x_2 \succ\dots$. 
\end{example}
%
\arka{drop either equivariance or well-order.
We make an educated guess and choose to drop the latter.
But even then reductions are non-terminating}

To make reductions terminating we disallow reductions to introduce new variables.
But now the definition of Gr\"{o}bener basis has to modified accordingly:
for a set $\Basis$ of polynomials to be a Gr\"{o}bner basis of an ideal $\idl[I]$,
for every $f\in\idl[I]$ there must be an element $b\in\Basis$ which does not use any variable that is not used by $f$,
and the leading monomial of $b$ divides the leading monomial of $f$\arka{itemise}.
In this article we call them \kl{strong Gr\"{o}bner bases}\arka{is this a good name?}.
We use ideas from \cite{GHOLAS24} to show that when $\monoid\actson\Indets$ is \kl{nicely ordered},
every $\monoid$-equivariant ideal has an orbit-finite \kl{strong Gr\"{o}bner base}.
However, now the naive extension of the Buchberger's algorithm to the equivariant setting does not compute \kl{strong Gr\"{o}bner bases}.
We illustrate this through concrete examples in \arka{write and refer}.
As one of the main technical contribution of this paper,
we give a simple modification of the Buchberger's algorithm that allows us to compute \kl{strong Gr\"{o}bner bases}.
%
\subsection{Organisation of the article}
%
\subsection{Related research}
%
\subsubsection{Equivariant ideals}
%
Equivariant polynomial ideals have received significant attention from the mathematics community in recent years. 
Above we have mentioned the articles that study existence and computability of orbit-finite Gr\"{o}bner bases of equivariant ideals.
In addition to that,
there also has been a lot of research on the structure of equivariant ideals and varieties.
Hilbert series of equivariant ideals have been studied by \cite{NaRo17,KrLeSn17,MaNa21,Na21}
The articles \cite{DEKL15,NaSn20} have studied equivariant varieties.
Prime and radical equivariant ideals have been studied by \cite{NaSn21,KuRi25,Sn24}.
Finally \cite{SaSn16,SaSn19,Yu25,NSY25} have studied equivariant varieties.
%
\subsubsection{Orbit-finitely generated algebraic structures}
%
Our results also fit into the research agenda of studying algebraic structures built from orbit-finite sets.
These algebraic objects have recently caught the attention of the theoretic computer science community due to their connection with models of computations with infinite alphabets,
such as register automata and data Petri nets
\cite{Boj13,BoSt20,Ste24,BFKM24,YBK26,GHL22,GHL25,Gho25,ramsey}.
To properly discuss this connection we need to set up some background material which is done in \arka{refer}.
The algebraic structures and their connections to automata theory has been discussed in \arka{refer}.
%
\subsection{Finite groups???????}
%
%In this setting one considers algebraic structures (monoids, vector spaces, polynomial rings) built from $\omega$-categorical relational structures $\Indets$.
%
%For instance, if $\Indets$ is the structure of \kl{equality atoms} (infinite set without any additional structure) $\aut{\Indets} = \symgr{\Indets}$.
%And if $\Indets$ is the structure of \kl{ordered atoms} (dense linear order without endpoints), $\aut{\Indets}$ is the set of all order preserving bijections of $\Indets$.
%By Ryll-Nardzewski theorem,
%$\Indets$ being $\omega$-categorical is equivalent to the action $\group\actson\Indets$ being oligomorphic
%(i.e., the $\Indets^n$ has finitely many $\group$-orbits for every $n$).
%Now one studies algebraic structures built from the relational structure $\Indets$.
%Mainly three classes of algebraic structures have been studied in the literature,
%one of which are polynomial rings of the form $\poly{\K}{\Indets}$.
%The other two being
%\begin{enumerate}
%\item monoids of endomorphisms of $\Indets^n$ \cite{Boj13,BoSt20,Ste24},
%\item freely generated vector spaces $\K(\Indets^n)$ \cite{BFKM24,YBK26,GHL22,GHL25,Gho25,ramsey}.
%\end{enumerate}
%
%
\tableofcontents
%
%\input{src/intro.tex}
%!TEX root = ../atomic_jsc.tex
% LTeX: language=en
\section{Preliminaries}
\label{sec:preliminaries}

%\subsection{Well-quasi-orders}
\arka{shift just before existence?}

\subsection{Partial orders, ordinals, well-founded sets, and well-quasi-ordered sets.}
\AP We assume basic familiarity with partial orders, well-founded sets, and ordinals.
We use the notation $\intro*\om$ for the first infinite ordinal (that is, $(\N, \leq)$), and write $X \intro*\ordplus Y$ for the lexicographic sum of two partial orders $X$ and $Y$
(i.e., $z < z'$ if either $z\in X$ and $z'\in Y$, or if $z < z'\in X$ or $z < z'\in Y$).
Similarly, the notation $X\intro*\ordtimes Y$ denotes the product of two partial orders equipped with the lexigographic ordering
(i.e., $(x_1, y_1) \leq (x_2, y_2)$ if either $x_1 < x_2$, or $x_1 = x_2$ and $y_1 \leq y_2$).
We also use the usual notations for finite ordinals, writing $\intro*\ordfin{n}$ for the finite ordinal of size $n$.
For instance, $\om \ordplus \ordfin{1}$ is the total order $\N \uplus\set{+\infty}$, where $+\infty$ is the new largest element.

\AP In order to guarantee the termination of the algorithms presented in this
paper, a key ingredient is the notion of \intro{well-quasi-ordering}
(WQO), which consists of sets $(X, \leq)$ such that every infinite sequence
$\seqof{x_i}[i \in \N]$ of elements of $X$ contains a pair $i < j$ such that
$x_i \leq x_j$. Examples of \kl{well-quasi-orderings} include finite sets with
any ordering, or $\N \times \N$ with the product ordering. We refer the reader
to \cite{SCSC17} for a comprehensive introduction to \kl{well-quasi-orderings}
and their applications in computer science.
\arka{Higman and Kruskal? Maybe in appendix?}
%
\subsection{Polynomials, monomials, divisibility.}
\arka{this subsection will be merged with equivariant grobner bases}
\AP We use the notation $\poly{\K}{\Indets}$ for the ring of polynomials with coefficients from a field $\K$ and indeterminates from a set $\Indets$,
and $\mon{\Indets}$ for the set of monomials in $\poly{\K}{\Indets}$.
Letters $p,q,r$ denote polynomials,
$\monelt,\monelt[n]$ denote monomials, and $a,b,\alpha,\beta$ denote coefficients in $\K$.
\arka{reserve one of them for variables or check instances where $a,b$ are used for variables}

A classical example of a \kl{WQO} is the set of monomials $\mon{\Indets}$,
endowed with the \kl{divisibility} relation $\intro*\divleq$ when $\Indets$
is finite. We recall that a monomial $\monelt[m]$ \intro{divides} a monomial
$\monelt[n]$ if there exists a monomial $\monelt[l]$ such that $\monelt[m]\monelt[l] = \monelt[n]$. In this case, we write $\monelt[m]
\reintro*\divleq \monelt[n]$.

\AP Let $\varleq$ be a total ordering on $\Indets$.
Using this order,
we define the \intro{co-lexicographic} (colex) ordering on monomials as follows:
$\monelt[n] \intro*\revlexlt \monelt[m]$ if there exists an
indeterminate $x \in \Indets$ such that $\monelt[n](x) < \monelt[m](x)$,
and such that for every $y \in \Indets$,
if $x \varlt y$ then $\monelt[n](y) = \monelt[m](y)$. Note that if $\monelt[n] \divleq \monelt[m]$,
then in particular $\monelt[n] \revlexleq \monelt[m]$. 

\AP We can now use the \kl{reverse lexicographic} ordering to identify particular elements in a given polynomial. Namely, for a polynomial $p \in \poly{\K}{\X}$, we define
the \intro{leading monomial} $\intro*\lm(p)$ of $p$ as the largest monomial
appearing in $p$ with respect to the \kl{revlex ordering}, and the
\intro{leading coefficient} $\intro*\lc(p)$ of $p$ as the coefficient of
$\lm(p)$ in $p$. We then define the \intro{leading term} $\intro*\lt(p)$ of
$p$ as the product of its \kl{leading monomial} and its \kl{leading
coefficient},
\arka{push to later}
and the \intro{characteristic monomial} $\intro*\cm(p)$ of $p$ as
the product of its \kl{leading monomial} and all the indeterminates appearing
in $p$.
We also define the \intro(monomial){domain} of $\monelt[m]$ as the set
$\intro*\dom(\monelt[m])$ of indeterminates $x \in \X$ such that $\monelt[m](x) \neq
0$. Because the coefficients and monomial in question are highly dependent on
the ordering $\varleq$, we allow ourselves to write $\lm[\Indets](p)$ to
highlight the precise ordered set of variables that was used to compute the
\kl{leading monomial} of $p$. We extend $\dom$ from 
monomials to polynomials by defining $\dom(p)$ as 
the union of the \kl(monomial){domains} of all monomials appearing in $p$.

\begin{remark}\label{rem:other-monomial-orderings}
In the case of a finite set of indeterminates, one can choose any total ordering
on $\mon{\Indets}$, as long as it contains the \kl{divisibility}
quasi-ordering, and is compatible with the product of monomials.\footnote{This
is often called a \emph{monomial ordering}, see \cite{CLO15}.} In our case,
having an infinite number of indeterminates, we rely on a connection between
$\lm(p)$ and $\dom(p)$: $\dom(p) \subseteq \dwset{\dom(\lm(p))}$, where
$\dwset{S}$ is the downwards closure of a set $S \subseteq \Indets$, i.e. the
set of all indeterminates $x \in \Indets$ such that $y \leq x$ for some $y \in
S$. This means that the \kl{leading monomial} encodes a \emph{global property}
of the polynomial, and it will be crucial in our termination arguments. This is
already at the core of the classical \emph{elimination theorems} \cite[Chapter 3, Theorem
2]{CLO15}.
\end{remark}

\subsection{Ideals and Gr\"{o}bner Bases}
\AP An \intro{ideal} $\idl$ of
$\poly{\K}{\X}$ is a non-empty subset of $\poly{\K}{\X}$ that is closed under
addition and multiplication by elements of $\poly{\K}{\X}$. Given a set $H
\subseteq \poly{\K}{\X}$, we denote by $\intro*\IdlGen{H}$ the ideal generated
by $H$, i.e. the smallest ideal that contains $H$. The \intro{ideal membership
problem} is the following decision problem: given a polynomial $p \in
\poly{\K}{\X}$ and a set of polynomials $H \subseteq \poly{\K}{\X}$, decide
whether $p$ belongs to the ideal $\IdlGen{H}$ generated by $H$. We know that
this problem is decidable when $\X$ is finite, and that it is even
$\EXPTIME$-complete \cite{MAME82}. The classical approach to the \kl{ideal
membership problem} is to use the \kl{Gröbner basis} theory, developed
in the 1970s by Buchberger~\cite{BUCH76}. 
A set $\Basis$ of polynomials is called a \intro{Gröbner basis} of
an ideal $\idl$ if, $\IdlGen{\Basis} = \idl$ and for every polynomial $p \in
\idl$, there exists a polynomial $q \in \Basis$ such that $\lm[\Indets](q)
\divleq \lm[\Indets](p)$.

Given a \kl{Gröbner basis} $\Basis$ of an ideal $\idl$, and a polynomial $p$,
it suffices to iteratively reduce the \kl{leading monomial} of $p$ by
subtracting multiples of elements in $\Basis$, until one cannot apply any
reductions. If the result is $0$, then $p$ belongs to $\idl$, and otherwise it
does not. 

\begin{example}
  Let $\Indets \defined \set{x,y,z}$ with $z < y < x$. The
  set $\Basis \defined \set{x^2y - z, x^2 - y}$ is not 
  a \kl{Gröbner basis} of the ideal $\idl$ it generates,
  because the polynomial $p \defined y^2 - z$ belongs to $\idl$
  but its \kl{leading monomial} $y^2$ is not divisible by 
  $\lm(x^2y - z) = x^2y$ nor by $\lm(x^2 - y) = x^2$.
\end{example}

%\para{Computability assumptions.}
%\AP We say that the action is
%\intro{effectively oligomorphic} if:
%%
%\begin{enumerate}
%\item It is \intro{oligomorphic}, i.e.\ for every $n \in \N$,
%$\Indets^n$ is \kl{orbit finite},
%\item There exists an algorithm that decides whether two elements $\vec{x},
%\vec{y} \in \Indets^*$ are in the same orbit under the action of 
%$\group$ on $\Indets^*$.
%\item There exists an algorithm which, on input $n\in\N$, outputs a set
%$A \subseteqfin \Indets^n$ such that $|A\cap U| = 1$ for every orbit
%$U \subseteq \Indets^n$.
%\end{enumerate}
%In particular, $\Indets$ itself is \kl{orbit finite} under the action of $\group$.
%
%
%A group action $\group \actson \X$ is said to be \intro(ord){compatible}
%with an ordering $\leq$ on $\X$ if for all $\gelem \in \group$ and $x,y \in
%\X$, we have $x \leq y$ if and only if $\gelem \cdot x \leq \gelem \cdot y$.
%Let us point out that in this case, $\revlexleq$ is also \kl(ord){compatible} with
%the action of $\group$ on $\mon{\X}$, i.e. for all $\gelem \in \group$ and
%monomials $\monelt[m], \monelt[n] \in \mon{\X}$, we have $\monelt[m] \revlexleq
%\monelt[n]$ if and only if $\gelem \cdot \monelt[m] \revlexleq \gelem \cdot
%\monelt[n]$.
%Our \intro{computability assumptions} on the tuple $(\Indets, \group,
%\leq)$ are that $\group$ acts \kl{effectively oligomorphic} on
%$\Indets$, and that its action is \kl(ord){compatible} with the ordering $\leq$
%on $\Indets$.
%
%\begin{example}
%  \label{ex:computability-assumptions}
%  Let $\Indets \defined \Q$ and $\group$ be the group of all
%  order preserving bijections of $\Q$.
%  Then, $\group$ acts \kl{effectively oligomorphically} on $\Indets$,
%  and its action is \kl(ord){compatible} with the ordering of $\Q$ by definition. 
%\end{example}
%
%Note that under our \kl{computability assumptions}, the set of polynomials
%$\poly{\K}{\Indets}$ is also \kl{effectively oligomorphic} under the action of
%$\group$ on $\Indets$ when restricted to polynomials with bounded degree.
%This is because a polynomial $p \in \poly{\K}{\Indets}$
%can be seen as an element of $(\K \times \Indets^{\leq d})^n$ where $n$ is the
%number of monomials in $p$, and $d$ is the maximal degree of a monomial
%appearing in $p$. Note that the set of all polynomials $\poly{\K}{\Indets}$ is
%not \kl{orbit finite}, precisely because the orbit of a polynomial $p$ under 
%the action of $\group$ cannot change the degree of $p$, and because there are 
%polynomials of arbitrarily large degree.
%
\subsection{Monoid actions, equivariance, orbit-finiteness and effective oligomorphism} 
\AP
In this section we define basic notions regarding monoid actions and define the concept of effective local oligomorphism.
As we show in \arka{refer to sec},
the latter condition will allow us to execute the equivariant Buchberger's algorithm (defined in \arka{refer to sec}).


A \intro{monoid} $\monoid$ is a set equipped with a binary operation that is associative, and has an identity element.
In our setting, we are interested in polynomial rings of the form $\poly{\K}{\Indets}$ where $\Indets$ is an infinite set of indeterminates that are equipped with an injective \intro{monoid action} $\monoid \actson \Indets$.
This means that there is monoid homomorphism $\psi_{\monoid}$ from $\monoid$ to the monoid of one to one endofunctions on $\Indets$.
Obviously, the same monoid $\monoid$ can act differently on a set $\Indets$.
However, usually the action $\psi_{\monoid}$ will be clear from the choice of $\monoid$ and $\Indets$ and we will use $\monoid\actson\Indets$ to specify this choice.
For each $\gelem \in \monoid$, we use $\gelem\cdot x$ to denote $\psi_{\monoid}(\gelem)(x)$.
Usually, we will be interested in actions $\monoid\actson\Indets$ where $\Indets$ comes with some additional structure and $\monoid$ is the monoid of all structure preserving injective endofunctions/group of all structure preserving bijections on $\Indets$.
For instance, $\inc{\OrderAtoms}\actson\OrderAtoms$ denote the obvious action of the order preserving one-to-one endofunctions on the structure $\OrderAtoms$ of rational numbers with order.
In \arka{cite sec} we give some more examples of such actions.  

\arka{shift to computability assumptions? Just after writing down theorem?
Computability follows from the rule of thumb that equivariant functions can be evaluated on orbit-finite sets.
We show how this rule of thumb formalises in our setting.
Maybe even formalise it?
Sets with atoms? Takes a lot of formalisms.
We take a shortcut. Input = a fixed numbers of (Tuple,string in binary). output similar.
But which kind of sets. maybe only lower levels in the hierarchy? bounded branching? heriditary finite sets or their orbits}
The action $\monoid\actson\Indets$ extends to actions $\monoid\actson\Indets^n$ for $n\in\N$ in the obvious way:
$\gelem\cdot(x_1,\dots,x_n) = (\gelem\cdot x_1,\dots,\gelem\cdot x_n)$ for $\gelem\in\monoid$ and $x_1,\dots,x_n\in\Indets$.
A subset $S \subseteq \Indets^n$ is called \intro{$\monoid$-equivariant} if for all $\gelem \in \monoid$ and $f \in S$, we have $\gelem \cdot f\in S$.
For example, $\setof{(x,y)\in\OrderAtoms^2}{x < y \in \OrderAtoms}$ is an $\inc{\OrderAtoms}$-equivariant subset of $\OrderAtoms^2$.
However, the subset $\setof{x \in \OrderAtoms}{x \text{ is an integer}}$ is not an $\inc{\OrderAtoms}$-equivariant subset of $\OrderAtoms$.

We denote $\intro*\orbit[\monoid]{P}$ for the set of all elements $\gelem \cdot p$ for $\gelem \in \monoid$ and $p \in P$.
When $P$ is a singleton set ${p}$,
we write $\orbit[\monoid]{p}$ instead of $\orbit[\monoid]{\{p\}}$.
An \kl{equivariant set} is said to be \intro{orbit finite} if it is the
union of finitely many \intro{orbits} under the action of $\monoid$.
Equivalently, an \reintro{orbit finite set}
is a set of the form $\orbit[\group]{P}$ for some finite set $P$.
%
\begin{definition}\label{def:loc oligo}
The action $\monoid\actson\Indets$ is called \intro{locally oligomorphic} if for every $x_1,\dots,x_n\in\Indets$,
the subset $\setof{(y_1,\dots,y_n)}{y_k\in\orbit[\monoid]{x_k}}$ is \kl{orbit-finite}.
The action $\monoid\actson\Indets$ is \intro{effectively locally oligomorphic} if:
\begin{enumerate}
\item it is \kl{locally oligomorphic},
\item there exists an algorithm to check whether a tuple $(a_1,\dots,a_n)\in\Indets^n$ is inside $\orbit[\monoid]{(b_1,\dots,b_n)}$ for some $(b_1,\dots,b_n)\in\Indets^n$, and
\item there exists an algorithm that takes elements $a_1,\dots,a_n\in\Indets$ as input and outputs a finite set $P$ such that $\setof{(y_1,\dots,y_n)}{y_k\in\orbit[\monoid]{x_k}} = \orbit[\group]{P}$.
\end{enumerate}
\end{definition}
%
\arka{maybe we need slightly different where tuples are chosen as input?}
%
The following lemma implies that \kl{local oligomorphicity} is a generalisation of the well-known notion of \kl{oligomorphicity} for group actions:
a group action $\group\actson\Indets$ is called \intro{oligomorphic} if sets of the form $\Indets^n$ is orbit-finite for every $n$.
This lemma will be followed by examples and non-examples of locally oligomorphic monoid actions. 
%
\begin{lemma}\label{lem:oligo loc oligo}
Let $\group\actson\Indets$ be a \kl{group action}.
\begin{enumerate}
\item If $\group\actson\Indets$ is \kl{oligomorphic} then it is \kl{locally oligomorphic}.
\item If $\Indets$ is \kl{orbit-finite} then $\group\actson\Indets$ is \kl{oligomorphic} if and only if it is \kl{locally oligomorphic}.
\end{enumerate}
\end{lemma}
%
\begin{lemma}\label{lem:orbit part}
Let $\group\actson\Indets$ be a group action.
For $(a_1,\dots,a_n),(b_1,\dots,b_n)\in\Indets^*$,
\[
(a_1,\dots,a_n)\in\orbit[\group]{(b_1,\dots,b_n)}
\quad\iff\quad
(b_1,\dots,b_n)\in\orbit[\group]{(a_1,\dots,a_n)}
\]
\end{lemma}
%
\arka{if this lemma is not used anywhere else, make it a claim and push it inside the proof of \Cref{lem:oligo loc oligo}.}

\arka{Maybe we can send the proof of \Cref{lem:oligo loc oligo} to the appendix?}
%
\begin{proof}
If $(a_1,\dots,a_n)\in\orbit[\group]{(b_1,\dots,b_n)}$ then there exists $\gelem\in\group$ such that
\[
(\gelem\cdot a_1,\dots,\gelem\cdot a_n) =
\gelem\cdot(a_1,\dots,a_n) =
(b_1,\dots,b_n) \ .
\]
This means $\gelem\cdot a_i = b_i$ for all $i$.
This implies
\[
\gelem^{-1}\cdot b_i = \gelem^{-1}\cdot(\gelem\cdot a_i) =
(\gelem^{-1}\gelem)\cdot a_i = \iden[\group]\cdot a_i  = a_i\ .
\]
Hence $\gelem^{-1}\cdot(b_1,\dots,b_n) = (a_1,\dots,a_n)$,
which means $(b_1,\dots,b_n)\in\orbit[\group]{(a_1,\dots,a_n)}$
\end{proof}
%
\begin{proof}[Proof of \Cref{lem:oligo loc oligo}]
\begin{enumerate}
\item \Cref{lem:orbit part} implies that for every $n$ the set $\Indets^n$ partitions into disjoint orbits.
Since $\group\actson\Indets$ is oligomorphic we know this partition is finite.
Let $U_1,\dots,U_N$ be the orbits inside $\Indets^n$.
Pick $x_1,\dots,x_n\in\Indets$.
The set $V = \setof{(y_1,\dots,y_n)\in\Indets^n}{y_k\in\orbit[\monoid]{x_k}}$ is an equivariant subset of $\Indets^n$.
Therefore for every orbit $U_{\ell}$ inside $\Indets^n$ either $U_{\ell} \subseteq V$ or $U_{\ell} \cap V = \emptyset$.
Therefore $V$ is also a finite union of orbits. 
\item The only if direction follows from the previous point.
We prove the if direction.
Let $A\subseteqfin\Indets$ be a finite set such that $\orbit[\group]{A} = \Indets$.
Then $\Indets^n$ is the union of sets of the form $K(a_1,\dots,a_n)\defined\setof{(b_1,\dots,b_n)}{b_i\in\orbit[\group]{a_i}}$ where $(a_1,\dots,a_n)\in A^n$.
By \Cref{lem:orbit part} the set $\Indets^n$ partitions into orbits.
Each orbit inside $\Indets^n$ must be inside $K(a_1,\dots,a_n)$ for some $(a_1,\dots,a_n)\in A^n$.
There are finitely many such sets.
Hence to prove $\Indets^n$ is orbit-finite it is enough to show that each $K(a_1,\dots,a_n)$ is orbit-finite.
This follows from local oligomorphicity of $\group\actson\Indets$.
\end{enumerate}
\end{proof}
%
Now we give some examples and non-examples of locally oligomorphic monoid actions.
%
\begin{example}
The action $\inc{\OrderAtoms}\actson\OrderAtoms$ is effectively locally oligomorphic.
We show that it satisfies the three points in \Cref{def:loc oligo}.
%
\begin{enumerate}
\item For every $a\in\OrderAtoms$,
we have $\orbit[\inc{\OrderAtoms}]{a} = \OrderAtoms$.
Therefore, for any subset $S\subseteq\OrderAtoms$ of size $n$,
we have $\setof{(b_1,\dots,b_n)}{b_k\in\orbit[\monoid]{a_k}} = \OrderAtoms^n = \orbit[\inc{\OrderAtoms}]{S^n}$.
This proves $\inc{\OrderAtoms}\actson\OrderAtoms$ is locally oligomorphic.
\item For every $(a_1,\dots,a_n),(b_1,\dots,b_n)\in\OrderAtoms^n$ we have
$\orbit[\monoid]{(a_1,\dots,a_n)} = \orbit[\monoid]{(b_1,\dots,b_n)}$ if and only if $(a_1,\dots,a_n)$ and $(b_1,\dots,b_n)$ have the same order-type (for every $i,j$ we have $a_i \leq a_j$ if and only if $b_i \leq b_j$),
which is easy to check.
\item Since $\setof{(b_1,\dots,b_n)}{b_k\in\orbit[\monoid]{a_k}} = \OrderAtoms^n = \orbit[\inc{\OrderAtoms}]{S^n}$ for every $(a_1,\dots,a_n)\in\OrderAtoms^n$ and $S\subseteqfin\OrderAtoms$ of size $n$,
on any input we can simply output $\{1,\dots,n\}^n$.
\end{enumerate}
The output of the algorithm in the last point contains redundancies since two tuples in $\{1,\dots,n\}^n$ can be in the same orbit (for example, $(1,1,2)$ and $(2,2,3)$).
However, the redundancies can be removed simply by using the algorithm that allows us to check if two tuples are in the same orbit.
\end{example}
%
\begin{example}
Let $\calZ$ denote the structure of integers with the usual ordering and $\aut{\calZ}$ the group of all order preserving bijections of $\calZ$.
Then $\aut{\calZ}\actson\calZ$ is not locally oligomorphic.
Observe that $\orbit[\aut{\calZ}]{a} = \calZ$ for every $a\in\calZ$.
Therefore by \Cref{lem:oligo loc oligo} we have to just show that $\aut{\calZ}\actson\calZ$ is not oligomorphic.
For every $a,b,c,d\in\calZ$ the tuples $(a,b)$ and $(c,d)$ are in the same orbit if and only if $a - b = c - d$.
Therefore the set $\calZ^2$ splits into infinitely many orbits,
making $\aut{\calZ}\actson\calZ$ not oligomorphic.
\end{example}
%
\section{Equivariant Gröbner bases}
%
\arka{merge with next section}
%
\subsection{Monoid actions}
%
A \intro{monoid} $\monoid$ is a set equipped with a binary operation that is associative, and has an identity element.
In our setting, we are interested in polynomial rings of the form $\poly{\K}{\Indets}$ where $\Indets$ is an infinite set of indeterminates that are equipped with an injective \intro{monoid action} $\monoid \actson \Indets$.
This means that there is monoid homomorphism $\psi_{\monoid}$ from $\monoid$ to the monoid of one to one endofunctions on $\Indets$.
Obviously, the same monoid $\monoid$ can act differently on a set $\Indets$.
However, usually the action $\psi_{\monoid}$ will be clear from the choice of $\monoid$ and $\Indets$ and we will use $\monoid\actson\Indets$ to specify this choice.
For each $\gelem \in \monoid$, we use $\gelem\cdot x$ to denote $\psi_{\monoid}(\gelem)(x)$.
Usually, we will be interested in actions $\monoid\actson\Indets$ where $\Indets$ comes with some additional structure and $\monoid$ is the monoid of all structure preserving injective endofunctions/group of all structure preserving bijections on $\Indets$.
For instance, $\inc{\OrderAtoms}\actson\OrderAtoms$ denote the obvious action of the order preserving one-to-one endofunctions on the structure $\OrderAtoms$ of rational numbers with order.
In \arka{cite sec} we give some more examples of such actions.
%
\subsection{Equivariant ideals}
%
We use the notation $\poly{\K}{\Indets}$ for the ring of polynomials with coefficients from a field $\K$ and indeterminates from a set $\Indets$.
The action $\monoid\actson\Indets$ extends to an action $\monoid\actson\poly{\K}{\Indets}$ in the obvious way.
For illustration, $\pi\cdot (x^2 + 2y) = (\pi\cdot x)^2 + 2(\pi\cdot y)$.
A subset $S \subseteq \poly{\K}{\Indets}$ is called \intro{$\monoid$-equivariant} if for all $\gelem \in \monoid$ and $f \in S$,
we have $\gelem \cdot f\in S$.
An \intro{$\monoid$-ideal} of $\poly{\K}{\Indets}$ is an ideal of $\poly{\K}{\Indets}$ that is $\monoid$-equivariant.
We give in \Cref{ex:idl-equiv} examples and non-examples of \kl{$\monoid$-ideals}.
%
\begin{example}\label{ex:idl-equiv}
\arka{Notational convension for ideal used here}
The set $\idl[S]_0 \subset \poly{\K}{\Indets}$ of all polynomials whose set of coefficients sums to $0$ is an $\symgr{\Indets}$-ideal.
The set $\idl[M]_x$ of all polynomials that are multiple of a fixed $x \in \Indets$ is an \kl{ideal} that is not a $\symgr{\Indets}$-ideal.
\end{example}
%
\begin{proof}
We leave it to the reader to verify that $\idl[S]_0$ is indeed a $\symgr{\Indets}$-ideal.

To see that $\idl[M]_x$ is not a $\symgr{\Indets}$-ideal pick any variable $y\neq x$.
Let $\pi_{x,y}\in\symgr{\Indets}$ be the permutation that swaps $x$ and $y$ acts as identity everywhere else.
Then $\pi_{x,y}\cdot x = y \notin \idl[M]_x$.
\end{proof}

%
%\begin{proof}
%    Let $p,q\in \idl[S]_0$, and $r \in \poly{\K}{\Indets}$.
%    Then, $p \times r + q$ is in $S_0$. Remark that 
%    $p,r$ and $q$ belong to a subset $\poly{\K}{\Indets}$ of the 
%    polynomials that uses only finitely many indeterminates.
%    In this subset, the sum of all coefficients is obtained
%    by applying the polynomials to the value $1$ for every indeterminate
%    $y \in \Indets$. We conclude that
%    $(p \times r + q)(1,\dots, 1) 
%    = p(1,\dots,1) \times r(1,\dots,1) + q(1,\dots,1)
%    = 0 \times r(1, \dots, 1) + 0 = 0$, hence that
%    $p \times r + q$ belongs to $S_0$. 
%    Because $0$ is in $S_0$, we conclude that $S_0$ is an \kl{ideal}.
%    Furthermore, if $\gelem \in \group$ and $p \in S_0$, then
%    the sum of the coefficients $\gelem \cdot p$ is exactly
%    the sum of the coefficients of $p$, hence is $0$ too.
%    This shows that $S_0$ is \kl(ideal){equivariant}.
%
%    It is clear that all multiples of a given polynomial $x \in \Indets$
%    is an ideal of $\poly{\K}{\Indets}$. This is not an \kl{equivariant ideal}:
%    take any bijection $\gelem \in \group$ that does not map $x$ to $x$ (it
%    exists because $\Indets$ is infinite and $\group$ is all permutations),
%    then $\gelem \cdot x$ is not a multiple of $x$, and therefore does 
%    not belong to the ideal.
%\end{proof}
\AP
%
The orbit of an element $p$ in $\poly{\K}{\Indets}$ is defined as
\[
\intro*\orbit[\monoid]{p} \defined
\setof{\gelem\cdot p}{\gelem\in\monoid} \ .
\]
For $P\subseteq \poly{\K}{\Indets}$ we use $\intro*\orbit[\monoid]{P}$ to denote the union of the orbits $\orbit[\monoid]{p}$ for elements in $p$:
\[
\orbit[\monoid]{P} = \bigcup_{p\in P}\orbit[\monoid]{p} \ .
\]
An \kl{$\monoid$-equivariant set} is said to be \intro{orbit finite} if it is the union of finitely many \intro{orbits} under the action of $\monoid$,
equivalently, it is a set of the form $\orbit[\group]{P}$ for some finite set $P$.
Obviously, not every \kl{equivariant subset} is \kl{orbit finite}.
For instance, $\setof{x^n}{x\in\Indets, n\in\N} \subseteq \poly{\K}{\Indets}$ is $\symgr{\Indets}$-equivariant but is orbit-finite since $x^n$ and $y^m$ can not be in the same orbit when $n=m$.

The ideal \intro*{$\IdlGen{G}$} generated by a subset $G\subseteq \poly{\K}{\Indets}$ is the smallest ideal containing $G$.
For instance, the ideals $S_0$ and $M_x$ in \Cref{ex:idl-equiv} are generated by the sets $\setof{y - 1}{y\in\Indets}$ and $\{x\}$, respectively.
An ideal is called \intro{orbit-finitely generated} if it is generated by an orbit-finite set.
For example, the generating set $\setof{y - 1}{y\in\Indets}$ of $S_0$ forms a single orbit under $\symgr{\Indets}$.
Since orbit-finite sets are equivariant by definition,
orbit-finitely generated ideals are also equivariant.
Cohen showed in \cite{Cohen67} that every $\symgr{\Indets}$-equivariant ideal in $\poly{\K}{\Indets}$ is orbit-finitely generated.
However, as the next example shows,
this statement does not hold in general for every monoid actions.
%
\begin{example}
Consider the action $\monoid_0\actson\Indets$ where $\monoid_0$ is the trivial monoid,
i.e., a monoid with identity as its only element.
Since $\monoid_0$-orbits are singletons,
any ideal in $\poly{\Q}{\Indets}$ is an $\monoid_0$-ideal
and an ideal is orbit-finitely generated if and only if it is finitely generated.
However, if $\Indets$ is infinite,
there are ideals in $\poly{\Q}{\Indets}$ that are not finitely generated.

For instance, consider the ideal $\idl[I]$ of all polynomials in $\poly{\Q}{\Indets}$ without any constant term.
Any finite subset $F$ of $\idl[I]$ can not generate $\idl[I]$ since every non-zero polynomial in $\IdlGen{F}$ must use some variable that is used by some polynomial in $F$.
Therefore, $y\in\idl[I]\setminus\IdlGen{F}$ for any variable $y$ that is not used by any polynomial in $F$.
\end{example}
%
\begin{example}
Consider the action $\symgr{A}\actson A^2$ where $A$ is an countably infinite set defined as: $\pi\cdot (a,b) = (\pi\cdot a, \pi\cdot b)$.
We give an example of an $\symgr{A}$-equivariant ideal in $\poly{\Q}{A^2}$ which is not orbit-finitely generated.

For $n\geq 3$, define
\[
C_n \defined \setof{(a_0,a_1)(a_1,a_2)\dots(a_{n-2},a_{n-1})(a_{n-1},a_0) \in \poly{\Q}{A^2}}{a_i\neq a_j\in A}
\]
The set $C_n$ is an $\symgr{A}$-orbit for every $n$.
Let $\calC$ be the $\symgr{A}$-ideal generated by $\cup_{n \geq 3} C_n$.
Then $\calC$ does not have any orbit-finite set of generators.

\arka{give intuitive reason with graphs}

\arka{orbit $\mapsto$ $\monoid$-orbit}

\arka{$\monoid$-equivariant ideal $\mapsto$ $\monoid$-ideal?}

\end{example}
%
\begin{proof}
Pick an arbitrary orbit-finite subset $F$ of $\calC$.
We show $\IdlGen{F} \subsetneq \calC$.
Since the group $\symgr{A}$ acts in a variable-wise manner on $\poly{\Q}{A^2}$,
every polynomial in a $\symgr{A}$-orbit inside $\poly{\Q}{A^2}$ uses the same number of variables.
Since $F$ is orbit-finite,
there exists $N\in\N$ such that every polynomial in $F$ uses at most $N$-variables.
Therefore, $F\subseteq \IdlGen{\cup_{n = 3}^N C_n}$, and as a consequence $\IdlGen{F}\subseteq \IdlGen{\cup_{n = 3}^N C_n}$.
We show that for any $m > N$ we have $C_m \subseteq \calC \setminus \IdlGen{\cup_{n = 3}^N C_n} 
\subseteq \calC \setminus \IdlGen{F}$.

\arka{notation for monomials used here}

Pick $\monelt[p] \in C_m$ for some $m > N$.
If $\monelt[p]\in \IdlGen{\cup_{n = 3}^N C_n}$ then
\[
\monelt[p] = \sum_{i = 1}^k r_i \monelt[q]_i\monelt[p]_i
\]
for some $r_i\in\Q$, $\monelt[q]_i\in\mon{A^2}$ and $\monelt[p]_i\in \cup_{n = 3}^N C_n$.
The equality can hold only if $\monelt[p] = \monelt[q]_i\monelt[p]_i$ for some $i$.
This implies $\monelt[p]_i$ divides $\monelt[p]$.
We show that this leads to a contradiction.

Let
\[
\monelt[p] = (a_0,a_1)(a_1,a_2)\dots(a_{m-2}a_{m-1})(a_{m-1},a_0)
\]
and
\[
\monelt[p]_i = (b_0,b_1)(b_1,b_2)\dots(b_{n-2}b_{n-1})(b_{n-1},b_0) \ .
\]
For positive integers $k$, Let $\oplus_k$ denote the addition of integers modulo $k$.
Since $\monelt[p]_i$ divides $\monelt[p]$, $(b_0,b_1) = (a_j,a_{j\oplus_m 1})$ for some $j$.
But then $(b_1,b_2) = (a_{j\oplus_m 1},a_{j\oplus_m 2})$,
By induction, we must have $(b_{0\oplus_n\ell},b_{1\oplus_n\ell}) = (a_{j\oplus \ell},a_{(j\oplus_m (\ell+1)})$ for every $\ell$.
The LHS of this equality takes at most $m$-values but the RHS takes $n$.
Since $n > N \geq m$, this leads to a contradiction.
\end{proof}
%
\arka{define the notation $\EqIdlGen{H}$}.


\AP A function $f \colon X \to Y$ between two sets $X$ and $Y$ equipped with
actions $\group \actson X$ and $\group \actson Y$ is said to be
\intro(func){equivariant} if for all $\gelem \in \group$ and $x \in X$, we have
$f(\gelem \cdot x) = \gelem \cdot f(x)$. For instance, the
\kl(monomial){domain} of a monomial is an \kl{equivariant function}: if $\gelem
\in \group$, then $\gelem \cdot \dom(\monelt[m]) = \dom(\gelem \cdot
\monelt[m])$. Let us point out that the image of an \kl{orbit finite set} under
an \kl{equivariant function} is \kl{orbit finite}, and that the algorithms that
we will develop in this paper are all \kl(func){equivariant}.
%
\subsection{Equivariant Gröbner bases}
%
\subsubsection{Divisibility and reduction}

\AP
We use the notation $\poly{\K}{\Indets}$ for the ring of polynomials with coefficients from a field $\K$ and indeterminates from a set $\Indets$,
and $\mon{\Indets}$ for the set of monomials in $\poly{\K}{\Indets}$.
Letters $p,q,r$ denote polynomials,
$\monelt,\monelt[n]$ denote monomials, and $a,b,\alpha,\beta$ denote coefficients in $\K$.
\arka{reserve one of them for variables or check instances where $a,b$ are used for variables}


We know from \cite{GHOLAS24} that a
necessary condition for the \kl{equivariant Hilbert basis property} to hold is
that the set  $\mon{\X}$  of monomials is a \kl{well-quasi-ordering} when
endowed with the \intro{divisibility up-to $\group$} relation
($\intro*\gdivleq$), which is defined as follows: for $\monelt_1, \monelt_2 \in
\mon{\X}$, we write $\monelt_1 \gdivleq \monelt_2$ if there exists $\gelem \in
\group$ such that $\monelt_1$ \kl{divides} $\gelem \cdot \monelt_2$. This
relation also extends to monomials that are functions from $\Indets$ to
$(Y,\leq)$ with finite support, where $Y$ is any partially ordered set. We say
that a set $\Basis \subseteq \poly{\K}{\X}$ is an \intro{equivariant Gröbner
basis} of an equivariant ideal $\idl$ if $\Basis$ is \kl{equivariant},
$\IdlGen{\Basis} = \idl$, and for every polynomial $p \in \idl$, there exists
$q \in \Basis$ such that $\lm[\Indets](q) \gdivleq \lm[\Indets](p)$ and
$\dom(q) \subseteq \dom(p)$, following the definition of \cite{GHOLAS24}.

\AP We say that a group action $\group \actson \Indets$ is
\intro{well-structured} if the set $(\mon[Y]{\Indets}, \gdivleq)$ is
\kl{well-quasi-ordered}, for every \kl{well-quasi-order} $(Y,\leq)$. We say
that it is \intro{$\omega$-well-structured} if $(\mon{\Indets}, \gdivleq)$ is
\kl{well-quasi-ordered}.

Note that even in the case of a finite set of variables, a \kl{Gröbner basis}
is not necessarily an \kl{equivariant Gröbner basis}, because of the
\kl(polynomial){domain} condition. However, every \kl{equivariant Gröbner
basis} is a \kl{Gröbner basis}.

\begin{example}
  \label{ex:equivariant-gb}
  Let $\Indets \defined \set{ x_1, x_2 }$,
  with $x_1 \varleq x_2$,
  and $\group$ be the trivial group.
  Let us furthermore consider the ideal $\idl$ \kl(idl){generated by}
  $\set{ x_1, x_2 }$.
  Then, the set $\Basis \defined \set{ x_2 - x_1, x_1 }$ is a
  \kl{Gröbner basis} of $\idl$, but not an \kl{equivariant Gröbner basis}.
  Indeed, $x_2 \in \idl$, but there is no polynomial $q \in \Basis$
  such that $\lm(q) \divleq x_2$ and $\dom(q) \subseteq \dom(x_2)$.
\end{example}

\subsection{Computability of equivariant Gr\"{o}bner bases}
%
\arka{detailed description of our results}

We denote $\intro*\orbit[\monoid]{P}$ for the set of all elements $\gelem \cdot p$ for $\gelem \in \monoid$ and $p \in P$.
When $E$ is a singleton set ${p}$,
we write $\orbit[\monoid]{p}$ instead of $\orbit[\monoid]{\{p\}}$.
An \kl{equivariant set} is said to be \intro{orbit finite} if it is the
union of finitely many \intro{orbits} under the action of $\monoid$.
Equivalently, an \reintro{orbit finite set}
is a set of the form $\orbit[\group]{P}$ for some finite set $P$.
%
%
%
%\arka{Does not satisfy the domain condition}
%In the finite case, one can always compute an \kl{equivariant Gröbner basis} by
%computing \kl{Gröbner bases} for every possible ordering of the indeterminates,
%and taking their union.\footnote{This algorithm is correct because we are
%  considering the \kl{reverse lexicographic} ordering.}
 





%\input{src/preliminaries.tex}
%!TEX root = ../atomic_jsc.tex
% LTeX: language=en
\section{Equivariant Gr\"{o}bner bases}
%
In this section we give the definition and study basic properties of equivariant Gr\"{o}bner bases of $\monoid$-equivariant ideals in $\poly{\K}{\Indets}$ where $\monoid\actson\Indets$ is a \kl{nicely ordered} monoid action.
To illustrate the definition and the properties,
we use the action $\inc{\OrderAtoms}\actson{\OrderAtoms}$.
\arka{use the same notation in the intro? maybe move it earlier}
Here, $\OrderAtoms$ denotes the set of rational numbers and $\inc{\OrderAtoms}$ denotes the set of strictly monotonic functions from $\OrderAtoms$ to itself.
We use the symbol $\OrderAtoms$ to emphasise the fact that here it is considered as a set with a linear order and not a field.
In the latter case we use the symbol $\Q$.
The action $\inc{\OrderAtoms}\actson{\OrderAtoms}$ is \kl{ordered} by definition,
but clearly not \kl{well ordered}.
It is however \kl{nicely ordered}.
For an intuitive argument,
$x^{k_1}_1\dots x^{k_m}_m\gdivleq[\inc{\OrderAtoms}]y^{\ell_1}_1\dots y^{\ell_n}_n$ for monomials $x^{k_1}_1\dots x^{k_m}_m$ and $y^{\ell_1}_1\dots y^{\ell_n}_n$ in $\poly{\K}{\OrderAtoms}$ with $x_1 < \dots < x_m \in\OrderAtoms$ and $y_1 < \dots y_n \in\OrderAtoms$
(i.e., there exists $\pi\in\inc{\OrderAtoms}$ such that $\pi(x_1)^{k_1}\dots \pi(x_m)^{k_m}$ divides $y^{\ell_1}_1\dots y^{\ell_n}_n$)
if and only if the sequence $k_1\dots k_m$ is smaller than the sequence $\ell_1\dots \ell_n$ in the subsequence ordering (i.e., there exists $1 \leq r_1 < \dots < r_m \leq n$ such that $k_i \leq \ell_{r_i}$ for all $i\in\{1,\dots,m\}$).
Since the subsequence ordering is a \kl{well-quasi-order},
$\gdivleq[\inc{\OrderAtoms}]$ is also so.


In \arka{cite section} we make this intuitive argument formal and show how it can be generalised to prove that some other monoid actions are also \kl{nicely ordered}.
%
\arka{push later}
\begin{lemma}
The relation $\gdivleq[\monoid]$ is a \kl{well-quasi-order} on $\mon{\Indets}$ if and only if there is an equivariant function $\varphi$ from $\mon{\Indets}$ to a well-quasi-order $(P,\leq)$ such that $\monelt[p] \gdivleq \monelt[q]$ if $\varphi(\monelt[p]) \leq \varphi(\monelt[q])$.
\end{lemma}
%
\begin{proof}
($\implies$)
Define a relation $\leq$ in $\setof{\orbit{\monelt[p]}}{\monelt[p]\in\mon{\Indets}}$ as
\[
\orbit{\monelt[p]} \leq \orbit{\monelt[q]}
\quad\text{when}\quad \monelt[p] \gdivleq[\monoid] \monelt[q]
\]
If $\gdivleq[\monoid]$ is a \kl{well-quasi-order} then so is $\leq$.
And the function $\varphi \defined \monelt[p] \mapsto \orbit{\monelt[p]}$ is an equivariant function from $\mon{\Indets}$ to the well-quasi-ordered set $(\setof{\orbit{\monelt[p]}}{\monelt[p]\in\mon{\Indets}},\leq)$ such that $\monelt[p] \gdivleq \monelt[q]$ if and only if $\varphi(\monelt[p]) \leq \varphi(\monelt[q])$.

\noindent($\impliedby$)
Pick an infinite sequence $\monelt[p]_1,\monelt[p]_2,\dots$ of monomials in $\mon{\Indets}$.
There must exist $i < j$ such that $\varphi(\monelt[p]_i) \leq \varphi(\monelt[p]_j)$.
But then $\monelt[p]_i \gdivleq \monelt[p]_j$.
\end{proof}
%
Let $\varleq$ denote the total order on $\Indets$ that is preserved by $\monoid\actson\Indets$.
\arka{$\mon{\Indets}$ hasn't been defined yet}
A monomial $\monelt[p]\in\mon{\Indets}$ is sometimes considered as a function from $\Indets$ to $\N$,
where $\monelt[p](x)$ is the exponent of $x$ in $\monelt[p]$.
For instance, if $\monelt[p] = x^2y$ then $\monelt[p](x) = 2$, $\monelt[p](y) = 1$ and $\monelt[p](z) = 0$ for all $z\neq x,y$.
Using $\varleq$, we define the \intro{co-lexicographic} (colex) ordering on monomials as
follows: $\monelt[n] \intro*\revlexlt \monelt[m]$ if there exists an
indeterminate $x \in \Indets$ such that $\monelt[n](x) < \monelt[m](x)$, and such
that for every $y \in \Indets$, if $x \varlt y$ then $\monelt[n](y) =
\monelt[m](y)$.
For example,
for $x \varlt y \varlt z \in\Indets$ we have
$x^2 y \revlexlt xy^2 \revlexlt z \revlexlt y^2z$.
%
Since $\varleq$ is not assumed to be a well-order and $x \varlt y$ implies $x \revlexlt y$ for $x,y\in\Indets$,
$\revlexleq$ can not be guaranteed to be a well-order.
However, it is still well-behaved with respect to multiplication.
\begin{lemma}
For every $\monelt[p],\monelt[q],\monelt[r]$ we have
\begin{enumerate}
\item $1 \revlexleq \monelt[p]$, and
\item if $\monelt[p] \revlexleq \monelt[q]$ then $\monelt[p]\monelt[r] \revlexleq \monelt[q]\monelt[r]$.
\end{enumerate}
\end{lemma}
%
In particular if $\monelt[n] \divleq \monelt[m]$ then $\monelt[n] \revlexleq \monelt[m]$. 

For a polynomial $p \in \poly{\K}{\X}$,
we define the \intro{leading monomial} $\intro*\lm(p)$ of $p$ as the largest monomial appearing in $p$ with respect to $\revlexleq$,
and the \intro{leading coefficient} $\intro*\lc(p)$ of $p$ as the coefficient of $\lm(p)$ in $p$.
We then define the \intro{leading term} $\intro*\lt(p)$ of $p$ as the product of its \kl{leading monomial} and its \kl{leading coefficient},
and the \intro{characteristic monomial} $\intro*\cm(p)$ of $p$ as the product of its \kl{leading monomial} and all the indeterminates appearing in $p$.
We also define the \intro(monomial){domain} of $\monelt[m]$ as the set $\intro*\dom(\monelt[m])$ of indeterminates $x \in \X$ such that $\monelt[m](x) \neq 0$.
Because the coefficients and monomial in question are highly dependent on the ordering $\varleq$,
we allow ourselves to write $\lm[\Indets](p)$ to highlight the precise ordered set of variables that was used to compute the \kl{leading monomial} of $p$.
\arka{is this needed?}
We extend $\dom$ from monomials to polynomials by defining $\dom(p)$ as the union of the \kl(monomial){domains} of all monomials appearing in $p$.
%
\begin{example}
Say $x \varlt y \varlt z \in \Indets$.
Let $p = x^2y^3 - \frac{1}{2}y^2z\in\poly{\Q}{\Indets}$.
Then $\lm(p) = y^2z$, $\lc(p) = -\frac{1}{2}$ and $\lt(p) = - \frac{1}{2}y^2z$.
\end{example}
%
\begin{remark}
The choice of the co-lexicographic ordering to order monomials is not completely arbitrary.
In \arka{refer to section} we discuss which properties of the co-lexicographic order becomes important for computing equivariant Gr\"{o}bner bases.
In this section we also show that some other well-known ordering of monomials such as the lexicographic ordering does not have these properties.
\end{remark}
%
\section{Equivariant Buchberger's algorithm}
%


\AP A function $f \colon X \to Y$ between two sets $X$ and $Y$ equipped with
actions $\group \actson X$ and $\group \actson Y$ is said to be
\intro(func){equivariant} if for all $\gelem \in \group$ and $x \in X$, we have
$f(\gelem \cdot x) = \gelem \cdot f(x)$. For instance, the
\kl(monomial){domain} of a monomial is an \kl{equivariant function}: if $\gelem
\in \group$, then $\gelem \cdot \dom(\monelt[m]) = \dom(\gelem \cdot
\monelt[m])$. Let us point out that the image of an \kl{orbit finite set} under
an \kl{equivariant function} is \kl{orbit finite}, and that the algorithms that
we will develop in this paper are all \kl(func){equivariant}.
%
\section{Weak Equivariant Gröbner Bases}
\label{sec:weakgb}

\AP In this section we prove that a natural adaptation of \kl{Buchberger's
algorithm} to the equivariant setting computes a \kl{weak equivariant Gröbner
basis} of an \kl{equivariant ideal}. This can be seen as an analysis of the
classical algorithm in the equivariant setting. We assume for the rest of
the section that $\Indets$ is a set of indeterminates equipped with a group
$\group$ acting \kl{effectively oligomorphically} on $\X$, and that $\X$ is
equipped with a total ordering $\varleq$ that is \kl(ord){compatible} with the
action of $\group$. The crucial object of this section is the notion of
\kl{decomposition} of a polynomial with respect to a set $H$.

\begin{definition}
  \label{def:decomposition}
  Let $H$ be a set of polynomials. A \intro{decomposition} of $p$
  with respect to $H$ is given by a finite sequence 
  $\mathfrak{d} \defined \seqof{(a_i, \monelt_i, h_i)}[i \in I]$ such that
  \begin{equation}
    p = \sum_{i \in I} a_i \monelt_i h_i \quad ,
  \end{equation}
  where $a_i \in \K$, $\monelt_i \in \mon{\X}$, and $h_i \in H$ for all $i \in I$.
  The \intro{domain of the decomposition}
  that we write $\intro*\domdec(\mathfrak{d})$ is defined as the union
  of the domains of the polynomials $\monelt_i h_i$ for all $i \in I$.
  The \intro{leading monomial of the decomposition} is defined as
  $
    \intro*\lmdec(\mathfrak{d}) \defined \max(\seqof{\lm(\monelt_i h_i)}[i \in I])
  $.
\end{definition}

Leveraging the notion of decomposition, we can define a weakening of the notion
of \kl{equivariant Gröbner basis}, that essentially mimics the classical notion
of \kl{equivariant Gröbner basis} at the level of \kl{decompositions} instead
of polynomials.

\begin{definition}
  An \kl{equivariant set} $\Basis$ of polynomials is 
  a \intro{weak equivariant Gröbner basis} of an \kl{equivariant ideal}
  $\idl$ if $\IdlGen{\Basis} = \idl$, and if for every polynomial $p \in \idl$,
  and decomposition $\mathfrak{d}$ of $p$ with respect to $\Basis$, there
  exists a decomposition $\mathfrak{d}'$ of $p$ with respect to $\Basis$ such that
  $\domdec(\mathfrak{d}') \subseteq \domdec(\mathfrak{d})$,
  and 
  such that $\lmdec(\mathfrak{d}') = \lm(p)$.
\end{definition}

\AP To compute \kl{weak equivariant Gröbner bases}, we use a rewriting
relation. Given $p,r \in \poly{\K}{\X}$, we write $p \intro*\toeucl{H}
r$ if and only if there exists $q \in H$, $a \in \K$, and $\monelt \in
\mon{\X}$ such that $p = a \monelt q + r$, $\dom(r) \subseteq \dom(p)$, and
$\lm[\Indets](r) \revlexlt \lm[\Indets](p)$. To simplify the
notations, we write $r \intro*\pmonlt p$ to denote $\dom(r) \subseteq
\dom(p)$, and $\lm[\Indets](r) \revlexlt \lm[\Indets](p)$, leaving the
ordered set of indeterminates $\Indets$ implicit.
The relation $\pmonleq$ is extended to decompositions by using 
the analogues of $\dom$ and $\lm$ for decompositions.

\begin{lemma}
  \label{lem:chm}
  The quasi-ordering $\pmonleq$ is \kl(ord){compatible} with the action of $\group$,
  and is well-founded on polynomials, and on \kl{decompositions} of polynomials.
\end{lemma}
\begin{proof}
  The first property is immediate because $\dom$, $\lm$, and $\revlexleq$ are
  compatible with the group action $\group$. 
  The second property follows from the fact that $\revlexlt$ is a total
  well-founded ordering whenever one has fixed finitely many possible 
  indeterminates. In a decreasing sequence, the support of the leading 
  monomials is also decreasing, so that sequence contains only finitely many 
  indeterminates, hence we conclude.
  The same proof works for decompositions.
\end{proof}

\AP As a consequence of \cref{lem:chm}, we know that the rewriting relation
$\toeucl{H}$ is \intro{terminating} for every set $H$. Given a set $H$ of
polynomials, and given a polynomial $p \in \poly{\K}{\X}$, we say that $p$ is
\intro{normalised} with respect to $H$ if there are no transitions $p
\toeucl{H} r$. The set of \intro{remainders} of $p$ with respect to $H$ is
denoted $\intro*\rem{H}{p}$, and is defined as the set of all polynomials $r$
such that $p \toeucl{H}^* r$ and $r$ is \kl{normalised} with respect to $H$.
The following lemma states that $\rem{H}{\cdot}$ is a computable function from
our setting.


\begin{lemma}[name={},restate={lem:normalisation}]
  \label{lem:normalisation}
  Let $H$ be an \kl{orbit finite set} of polynomials, and let $p \in \poly{\K}{\X}$ be a
  polynomial. Then $\rem{H}{p}$ is finite.
  Furthermore, this computation
  is \kl(func){equivariant}. In particular, 
  $\rem{H}{K}$ is a computable \kl{orbit finite} set 
  for every \kl{orbit finite} set $K$ of polynomials.
  \proofref{lem:normalisation}
\end{lemma}


\AP Now that we have a quasi-ordering on polynomials, we prove that given
an \kl{orbit finite} set $H$ of generators, we can compute a \kl{weak
equivariant Gröbner basis}. The computation closely follows the classical
\kl{Buchberger's algorithm}. The main idea is to saturate the set of
generators $H$ to remove some \emph{critical pairs} of the rewriting relation
$\toeucl{H}$. Namely, given two polynomials $p$ and $q$ in $H$, we compute the
set $\intro*\CancelPoly{p}{q}$ of cancellations between $p$ and $q$ as the set of
polynomials of the form $r = \alpha \monelt[n] p + \beta \monelt[m] q$ such
that $\lm(r) < \max(\monelt[n] \lm(p), \monelt[m]\lm(q))$, where $\alpha,\beta
\in \K$, and where $\monelt[n], \monelt[m] \in \mon{\X}$. Let us recall that
given two monomials $\monelt[n], \monelt[m] \in \mon{\X}$, one can compute
$\intro*\lcm(\monelt[n], \monelt[m])$ as the least common multiple of the two
monomials, and that this is an \kl{equivariant operation}.
Using this, we can introduce the \intro{S-polynomial} of two polynomials $p$ and $q$
as in \cref{eq:spoly}.
 \begin{equation}
    \label{eq:spoly}
    \intro*\spoly{p}{q} \defined
    \frac{\lcm(\lm(p), \lm(q))}{\lt(p)} \times p
    - \frac{\lcm(\lm(p), \lm(q))}{\lt(q)} \times q
    \quad .
  \end{equation}



\begin{lemma}[restate=lem:spoly,name={S-polynomials}]
  \label{lem:spoly}
  Let $p$ and $q$ be two polynomials in $\poly{\K}{\X}$.
  All the polynomials in $\CancelPoly{p}{q}$ are obtained by multiplying a monomial
  with
  their \kl{S-polynomial} $\spoly{p}{q}$.
  \proofref{lem:spoly}
\end{lemma}


Remark that the \kl{S-polynomial} is equivariant: if $\gelem \in \group$, then
$\spoly{\gelem \cdot p}{\gelem \cdot q} = \gelem \cdot \spoly{p}{q}$. Given a
set $H$, we write $\intro*\spolyset(H) \defined \bigcup_{p,q \in H}
\rem{H}{\spoly{p}{q}}$. We are now ready to define the saturation algorithm
that computes \kl{weak equivariant Gröbner bases}, described in
\cref{alg:weakgb}.
Let us remark that \cref{alg:weakgb}
is an actual algorithm (\cref{lem:weakgb-computable}) that is
\kl(func){equivariant}.

\begin{algorithm}
  \caption{Computing \kl{weak equivariant Gröbner bases} using the algorithm \intro{weakgb}.}
    \label{alg:weakgb}
    \KwIn{An orbit finite set $H$ of polynomials}
    \KwOut{An orbit finite set $\Basis$ that is a \kl{weak equivariant Gröbner basis} of
      $\EqIdlGen{H}$}
    \Begin{
        $\Basis \gets H$\;
        \Repeat{$\Basis$ stabilizes}{
            $\Basis \gets \Basis \cup \spolyset(\Basis)$\;
        }
        \Return{$\Basis$}\;
    }
\end{algorithm}
%
\begin{lemma}
  \label{lem:weakgb-computable}
  \cref{alg:weakgb} is computable and \kl(func){equivariant},
  and produces an \kl{orbit finite} set $\Basis$ if it terminates.
\end{lemma}
%
\begin{proof}
  Observe that it is enough to show that $\spolyset{\Basis}$ is orbit-finite for every orbit-finite set $\Basis$.
  First, we compute $\Basis^2$, which is an \kl{orbit finite} set of pairs,
  because $\Basis$ is \kl{orbit finite} and $\Indets$ is 
  \kl{effectively oligomorphic}.
  Then, noting that $\spoly{-\ }{-}$ is computable and \kl(func){equivariant},
  we conclude
  that $\bigcup_{p,q \in H}{\spoly{p}{q}}$
  is computable and orbit-finite.
  Now, using \Cref{lem:normalisation}, one can compute the set $\spolyset(\Basis)$, which is also orbit-finite.
  Furthermore, one can decide whether the set $\Basis$ stabilises,
  because the membership of a polynomial $p$ in $\Basis$ is decidable,
  since $\group\actson\Indets$ is \kl{effectively oligomorphic} and $\Basis$ is \kl{orbit finite}.
\end{proof}


Let us now use the semantic assumptions to prove the termination of
\cref{alg:weakgb}
(\cref{lem:weakgb-termination})
and the correctness of the resulting orbit finite set
(\cref{lem:weakgb-correctness}).

\begin{lemma}[name={},restate={lem:weakgb-termination}]
  \label{lem:weakgb-termination}
  Assume that the action $\group \actson \Indets$ is \kl{$\omega$-well-structured}.
  Then, 
  \cref{alg:weakgb} terminates
  on every \kl{orbit finite} set $H$ of polynomials.
  \proofref{lem:weakgb-termination}
\end{lemma}


\begin{lemma}
  \label{lem:weakgb-correctness}
  Assume that $\Basis$ is the output of \cref{alg:weakgb}. Then, it 
  is a \kl{weak equivariant Gröbner basis} of the ideal
  $\EqIdlGen{H}$.
\end{lemma}
\begin{proof}
  It is clear that $\Basis$ is a generating set of $\EqIdlGen{H}$, because
  one only add polynomials that are in the ideal generated by $H$ at every step.

  Let $p \in \EqIdlGen{H}$ be a polynomial,
  and let $\mathfrak{d}$ be a \kl{decomposition} of $p$ with respect to
  $\Basis$, that is, a \kl{decomposition} of the form
  \begin{equation}
    p = \sum_{i \in I} \alpha_i \monelt_i p_i
    \quad .
  \end{equation}
  Where $\alpha_i \in \K$, $p_i \in \Basis$, and $\monelt_i \in \mon{\X}$,
  for all $i \in I$.

  Leveraging \cref{lem:chm}, we know that the ordering $\pmonleq$ is
  well-founded. As a consequence, we can consider a minimal 
  decomposition $\mathfrak{d}'$ of $p$ with respect to $\Basis$ such that $\mathfrak{d}'
  \pmonleq \mathfrak{d}$. We now distinguish two cases, depending on whether
  the leading monomial $\lmdec(\mathfrak{d}')$ of the decomposition $\mathfrak{d}'$ is
  equal to the leading monomial of $p$ or not.

  \begin{description}
    \item[Case 1:] $\lmdec(\mathfrak{d}') = \lm(p)$.
      In this case, we conclude immediately,
      as we also have by assumption $\dom(\mathfrak{d}') \subseteq \dom(\mathfrak{d})$.

    \item[Case 2:] $\lmdec(\mathfrak{d}') \neq \lm(p)$.
      In this case, it must be that the set $J$ of indices such that
      $\monelt[l] \defined \lm(\monelt_i p_i) = \lmdec(\mathfrak{d}')$ is non-empty.
      Let us remark that 
      the sum of leading coefficients 
      of the polynomials in $J$ must vanish: $\sum_{i \in J} \alpha_i \lc(p_i) = 0$.
      As a consequence, the set $J$ has size at least $2$.
      Let us distinguish one element $\star \in J$, and 
      write $J_\star = J \setminus \set{\star}$.
      We conclude that 
      $\alpha_\star = - \sum_{i \in J_\star} \alpha_i \lc(p_i) / \lc(p_\star)$.
      Let us now rewrite $p$ as follows:
      \begin{equation}
        p = \sum_{i \in J_\star} \alpha_i 
        \left(\monelt_i p_i - \frac{\lc(p_i)}{\lc(p_\star)} \monelt_\star p_\star\right)
        + \sum_{i \in I \setminus J} \alpha_i \monelt_i p_i
        \quad .
      \end{equation}
      Now, by definition,
      polynomials $\alpha_i \monelt_i p_i$ for $i \in I \setminus J$ have 
      leading monomials
      strictly smaller than $\monelt[l]$.
      Furthermore,
      the polynomials
      $\monelt_i p_i - \frac{\lc(p_i)}{\lc(p_\star)} \monelt_\star p_\star$ for $i \in J_\star$
      cancel their leading monomials, hence they belong
      to the set $\CancelPoly{p_i}{p_\star}$.
      By \cref{lem:spoly}, we know that these polynomials are obtained by
      multiplying the \kl{S-polynomial} $\spoly{p_i}{p_\star}$ by some monomial.
      Because \cref{alg:weakgb} terminated, we know that 
      $\spoly{p_i}{p_\star} \toeucl{\Basis}^* 0$ by construction.

      By definition of the rewriting relation, we conclude that one can rewrite
      $\spoly{p_i}{p_\star}$ as combination of polynomials in $\Basis$ that
      have smaller or equal leading monomials, and do not introduce new
      indeterminates.

      We conclude that
      the whole sum is composed of polynomials with leading monomials 
      strictly smaller than $\monelt[l]$, and using a subset of the indeterminates
      used in $\mathfrak{d}'$, leading to a contradiction
      because of the minimality of the latter. 
      \qedhere
  \end{description}
\end{proof}

As a consequence of the above lemmas, we can now conclude that the 
\cref{alg:weakgb} computes a \kl{weak equivariant Gröbner basis} of the
ideal $\EqIdlGen{H}$, as stated in \cref{thm:weakgb-comput}.

\begin{theorem}
  \label{thm:weakgb-comput}
  Assume that the action $\group \actson \Indets$
  is \kl{$\omega$-well-structured}
  and satisfies our \kl{computability assumptions}.
  Then, the algorithm $\mathsf{weakgb}$ that takes as input an \kl{orbit finite} set $H$ of generators of an
  \kl{equivariant ideal} $\idl$ and computes a \kl{weak equivariant Gröbner
  basis} $\Basis$ of $\idl$.
\end{theorem}

%\input{src/weakgb.tex}
\input{src/fullbasis.tex}
\input{src/examples.tex}
\input{src/applications.tex}
\input{src/undecidability.tex}
\input{src/conclusion.tex}

% Include the bibliography
\bibliographystyle{plainurl}
\bibliography{papers.bib}

% If there are any appendices, we include them here.
\appendix
\input{src/appendix.tex}


\end{document}
